\chapter{Question tables}\label{app:questions}
% STATUS: written 2026-11-06 (every row names a problem stated with its located
% citation in the section that asks its question; no new claim is made here. The
% question families are the eight fixed in
% Section~\ref{sec:basic-theory:question-families}.)
%
% For each question family and class of algebra, whether the answer is known,
% false with a counterexample, or open, with the source.

Section~\ref{sec:basic-theory:question-families} fixed eight questions that the
book asks of every class of algebra.
The chapters answer them one class at a time.
This appendix collects the answers in one place, family by family.
An answer is one of three: known, with the theorem that settles it; false, with
a counterexample; or open, with the source that states it as a question.
Each open row cites the problem number, and the problem is stated in full in
the section that asks its question.

% Freudenburg 2017 (pp. 81, 107--112, 281--293), Bavula 2017 (pp. 1--2),
% Belov-Kanel--Kontsevich 2007 (pp. 1--2), Hadwin--Li 2004 (pp. 703--713),
% Bresar--Villena 2002 (pp. 710--720), Martindale 1964 (pp. 183--187),
% Ayupov--Yusupov 2020 (pp. 1--8). Excerpts in
% ledger/projects/derivations-encyclopaedia/excerpts/.
This chapter collects the questions that no chapter of the book could answer.
Section~\ref{sec:basic-theory:question-families} fixed eight questions for the
book to ask of every class of algebra. % claim-ok: our own description of the structure of the book
It asked for one of three answers: known, false with a counterexample, or open.
The open answers are collected here, so that a reader can see the frontier of
the subject in one place.

\hbfigure{22-open-problems/question-families}{The question families that organise the open problems.}{fig:open:families}

Most of these problems have a named author, and the chapters met them one
at a time.
The Zariski--Lipman conjecture carries the names of Zariski and of Lipman, and
Lipman proved that a free module of derivations forces the local ring to be
normal \cite[p.~270]{BiswasGurjarKolte2014}.
The Jacobian conjecture stands behind the theory of locally nilpotent
derivations, and it can be stated in the language of derivations
\cite[p.~81]{Freudenburg2017}.
The conjecture of Dixmier has been open since 1968 for every dimension
\cite[pp.~1--2]{Bavula2017}.
Larson asks which algebras have a reflexive space of derivations
\cite[p.~703]{HadwinLi2004}.
\v{S}emrl introduced 2-local derivations in 1997 and asked when such a map is a
derivation \cite[p.~1]{AyupovYusupov2020}.
Singer and Wermer proved their radical theorem in 1955 and at the same time
stated the conjecture that continuity is not needed
\cite[p.~710]{BresarVillena2002}.
Martindale studied the Lie derivations of a primitive ring in 1964 and said that
his result falls far short of the conjecture of Herstein
\cite[p.~183]{Martindale1964}.

Each of those answers holds for the moment the source states it.
A problem can move from open to known over decades.

\begin{example}[{\cite[p.~710]{BresarVillena2002}}]\label{ex:appendix:three-answers}
Singer and Wermer proved in 1955 that a continuous derivation on a commutative
Banach algebra takes its values in the radical.
At the same time they stated the conjecture that the continuity hypothesis is
superfluous.
Johnson proved that conjecture for semisimple algebras in 1969.
Thomas settled it for every commutative Banach algebra more than thirty years
later \cite[p.~710]{BresarVillena2002}.
The noncommutative form of the same question is still open, and
Problem~\ref{prob:open:sw} states it.
\end{example}

Each open problem is stated in the chapter that asks it, with the page where
the literature makes it. % claim-ok: a statement about how the book is arranged
This appendix gathers them instead of restating them: every row names the
problem and the chapter that carries the discussion.
A problem that a chapter settles for one class of algebras is recorded in
Section~\ref{sec:basic-theory:open-problems-by-class} under that class.
The list is a snapshot of the book as it now stands.

The problems the chapters state, in the order that
Section~\ref{sec:basic-theory:question-families} fixed: the
Zariski--Lipman problem, the Jacobian and Dixmier problems, and the questions of
polynomial rings about finite generation; then the question of Larson on reflexivity,
the 2-local problem, the noncommutative Singer--Wermer conjecture, the Lie
problem of Herstein, and the four problems the chapters of the book raise.
Section~\ref{sec:basic-theory:open-problems-by-class} sorts the same list by
where the algebra lives: polynomial algebras, finite-dimensional algebras,
$C^{*}$-algebras, Lie algebras, prime rings and triangular algebras.
For each class it records which problem is settled there and which is not.

\section{Family 1: which kind implies which}

The book records the implications between kinds of derivation-like map.
Table~\ref{tab:q:implications} gives the ones the chapters settle.

\begin{table}[htbp]
  \centering
  \caption{Answers for the family ``which kind of derivation implies which''.}
  \label{tab:q:implications}
  \begin{tabular}{p{0.30\textwidth}p{0.26\textwidth}p{0.30\textwidth}}
    \toprule
    Class of algebra & Answer & Source \\
    \midrule
    Triangular algebra & $f$-derivations are not all determined; the
      question of which $f$ admit only inner ones is open &
      Problem~\ref{prob:open:fderiv} \\
    Prime ring, characteristic not $2$ & symmetric biderivations: the
      expected description fails without a further hypothesis & open,
      Problem~\ref{prob:open:vukman} \\
    Semiprime ring & additive maps with $[f(x),x]_n=0$: the commuting
      case is settled, the general one is not & open,
      Problem~\ref{prob:open:commuting} \\
    \bottomrule
  \end{tabular}
\end{table}

\section{Family 2: when two definitions are equivalent}

Section~\ref{sec:basic-theory:family-tree} named the kinds that carry more than
one definition in the literature.
Where the definitions agree, a chapter says so and cites both sources; where
they do not, the row below records it.

\begin{table}[htbp]
  \centering
  \caption{Answers for the family ``when are two definitions equivalent''.}
  \label{tab:q:definitions}
  \begin{tabular}{p{0.30\textwidth}p{0.26\textwidth}p{0.30\textwidth}}
    \toprule
    Class of algebra & Answer & Source \\
    \midrule
    Rings with an $f$ & three notions carry the name $f$-derivation; the
      one used here is fixed by a defining identity & the census of
      Section~\ref{sec:basic-theory:family-tree} \\
    \bottomrule
  \end{tabular}
\end{table}

\section{Family 3: is every derivation inner}

This is the question of $H^1$, and the book answers it class by class.

\begin{table}[htbp]
  \centering
  \caption{Answers for the family ``is every derivation inner''.}
  \label{tab:q:inner}
  \begin{tabular}{p{0.30\textwidth}p{0.26\textwidth}p{0.30\textwidth}}
    \toprule
    Class of algebra & Answer & Source \\
    \midrule
    Central simple algebra & known: every derivation into a bimodule is
      inner & Theorem~\ref{thm:basic:csa-inner} \\
    Simple ring & Lie derivations: whether each is a derivation plus a map
      into the centre is open & Problem~\ref{prob:open:herstein} \\
    Local ring, characteristic $0$ & known for a free module of derivations;
      the general case is the Zariski--Lipman conjecture & open,
      Problem~\ref{prob:open:zl} \\
    Projective surface & the logarithmic case is open &
      Problem~\ref{prob:open:zl-projective} \\
    \bottomrule
  \end{tabular}
\end{table}

\section{Family 4: is every derivation continuous}

Continuity is the question of this family, and it is settled for some classes
and open for others.

\begin{table}[htbp]
  \centering
  \caption{Answers for the family ``is every derivation continuous''.}
  \label{tab:q:continuity}
  \begin{tabular}{p{0.30\textwidth}p{0.26\textwidth}p{0.30\textwidth}}
    \toprule
    Class of algebra & Answer & Source \\
    \midrule
    Banach algebra & whether a derivation with no continuity assumed leaves
      every primitive ideal invariant is open & Problem~\ref{prob:open:sw} \\
    Algebra of bounded operators & whether $\operatorname{der}(A)$ is
      reflexive is open for general $A$ & Problem~\ref{prob:open:larson} \\
    Lie algebra, characteristic $0$ & whether every 2-local derivation is a
      derivation is open in infinite dimension &
      Problem~\ref{prob:open:twolocal} \\
    \bottomrule
  \end{tabular}
\end{table}

\section{Family 5: does the algebra have to be commutative}

the theorem of Posner and the functional identities answer this for prime rings.
The noncommutative cases that remain open are collected here.

\begin{table}[htbp]
  \centering
  \caption{Answers for the family ``does the algebra have to be
  commutative''.}
  \label{tab:q:commutative}
  \begin{tabular}{p{0.30\textwidth}p{0.26\textwidth}p{0.30\textwidth}}
    \toprule
    Class of algebra & Answer & Source \\
    \midrule
    Prime ring & known: a nonzero derivation $d$ with $ad(a)-d(a)a$ central
      for every $a$ forces the ring to be commutative &
      Theorem~\ref{thm:prime:posner-two} \\
    $C^{*}$-algebra & Jordan derivations into a bimodule: the general case is
      open & Problem~\ref{prob:open:jordan-cstar} \\
    \bottomrule
  \end{tabular}
\end{table}

\section{Family 6: can all derivations be described}

A class whose derivations can be listed has a structure theorem.
The classes where the list is still unknown carry the open rows.

\begin{table}[htbp]
  \centering
  \caption{Answers for the family ``can all derivations of the class be
  described''.}
  \label{tab:q:describe}
  \begin{tabular}{p{0.30\textwidth}p{0.26\textwidth}p{0.30\textwidth}}
    \toprule
    Class of algebra & Answer & Source \\
    \midrule
    Polynomial ring & locally nilpotent derivations: the homogeneous case
      with $D^2x_i=0$ is open beyond the known dimensions &
      Problem~\ref{prob:open:hdp} \\
    Weyl algebra & whether every endomorphism is an automorphism, the
      Dixmier conjecture, is open since 1968 &
      Problem~\ref{prob:open:dixmier} \\
    Polynomial map, characteristic $0$ & the Jacobian conjecture is open &
      Problem~\ref{prob:open:jc} \\
    \bottomrule
  \end{tabular}
\end{table}

\section{Family 7: does a derivation extend or lift}

Section~\ref{sec:basic-theory:extension} fixed this question.
The rows below record where the extension question is still open.

\begin{table}[htbp]
  \centering
  \caption{Answers for the family ``does a derivation extend or lift to a
  larger algebra''.}
  \label{tab:q:extension}
  \begin{tabular}{p{0.30\textwidth}p{0.26\textwidth}p{0.30\textwidth}}
    \toprule
    Class of algebra & Answer & Source \\
    \midrule
    Unipotent group acting on affine space & whether the action can be
      triangularized, the question of Bass, is open &
      Problem~\ref{prob:open:bass} \\
    \bottomrule
  \end{tabular}
\end{table}

\section{Family 8: where does an expected implication fail}

A counterexample answers this family.
The counterexamples the book records are in the chapters named below, and the
one open case is listed last.

\begin{table}[htbp]
  \centering
  \caption{Answers for the family ``where does an expected implication
  fail''.}
  \label{tab:q:counterexamples}
  \begin{tabular}{p{0.30\textwidth}p{0.26\textwidth}p{0.30\textwidth}}
    \toprule
    Class of algebra & Answer & Source \\
    \midrule
    Prime ring, prime characteristic & false: the Zariski--Lipman
      statement fails & Chapter~\ref{ch:commutative} \\
    Simple Lie algebra, characteristic $0$ & false: outer derivations exist
      where the characteristic-zero answer predicts none &
      Chapter~\ref{ch:lie-algebras} \\
    \bottomrule
  \end{tabular}
\end{table}

\section{Reading the tables}

Every row above names a chapter or a problem number, and each of those carries
the located citation in its own place.
No row is a summary of a summary: a reader who wants the statement turns to the
chapter or the problem and finds it there, with its page.
The tables are an index to those statements, and they are refreshed with each
edition as the chapters change.
